\scanpage{168}
\mspage{42}

\unsure{grossièrement|modérément} saturé dans $\unsure{\mathcal M|M}$, stable par $\unsure{\underline{\lim}|\underrightarrow{\lim}}$ filtrantes,
$C_0$ une \uline{petite} partie de $\mathrm{Fl}\mathcal M$, \deleted{\unsure{contenue}} telle
que
\[
 C_0\subset\mathrm{Cof}_W,\qquad \underset{\substack{\Vert\\(TF)}}{C_{0*}}\subset W,
\]
\unsure{posons}
\[
 C=\widetilde C_0,\ TC=C\cap W,\ TF=C_{\unsure{*}},\ F=(TC)_*,
\]
et soit $TC_0$ \deleted{une petite} partie de $TC=C\cap W$.
On a alors \deleted{$\unsure{TC_0\subset TC}$, $\unsure{F_0\supset F}$}
\deleted{Dans ces conditions \unsure{pour tout} $C_0$, il y a équivalence :}
% margin note written sideways
\margin{$(4.4.6)\ \left\{\begin{aligned}
&C\subset\mathrm{Cof}_W,\ \deleted{\unsure{TC_0\subset C}}\ TC=C\cap W^{\mathrm{univ}}\quad\text{(pour mémoire)}\\
&\widetilde{TC_0}\subset TC,\quad F_0\supset F\quad\text{(\unsure{ce dernier|ce dernier est} \deleted{\unsure{lieu}} trivial)}
\end{aligned}\right.$}
\deleted{$\unsure{(4.4.6)}\quad\unsure{TC_0}\quad\widetilde{TC_0}\subset TC,\quad\unsure{(TC_0)_*}\ F_0\overset{\mathrm{déf}}{=}(TC_0)_*\supset F$,}
et les conditions suivantes sont équivalentes,
\[
\begin{aligned}
\mathrm{a})&\quad \widetilde{TC_0}=TC\ \deleted{\unsure{\subset}} \quad\text{(on a $\subset$)}\\
\mathrm{a}')&\quad \deleted{\unsure{TC_*=(TC_0)_*}}\ F_0=F\quad\text{(on a $\supset$)}\\
\mathrm{b})&\quad TF=F_0\cap W\quad\text{(on a $\subset$)}\\
\mathrm{c})&\quad \unsure{a')|c')}W=TF\circ\widetilde{TC_0}\quad\text{(on a $\supset$)}
\end{aligned}
\]
\margin{\uline{NB} On voit facilement que $\boxed{TF=F\cap W}$ mais \unsure{ce|ceci} ne suffit pas pour avoir $TF=F_0\cap W$ pour $TC_0$ \unsure{donné}.}
\margin{\uline{NB} On sait que $\boxed{W=TF\circ TC}$, ce qui n'est nullement suffisant pour impliquer c) !}
\deleted{[\unsure{d)}\inserted{\unsure{bis}}\ \unsure{$(W,C,TC)$ satisfait à factorisation}}
et impliquent que $(W,C,F)$ est un triple
\uline{de Quillen \unsure{clos}.}

\uline{Dém} L'inclusion $\widetilde{TC_0}\subset TC$ résulte de la Q-stabilité
% slanted
\inserted{à g. de $TC$ (\unsure{cor.} 4.4.4)
et du fait que $TC_0\subset TC$ est petite
(\unsure{à spécifier|à figurer} dans le th. !
tout comme pour $C_0\subset\mathrm{Cof}_W$ \ldots)}
\deleted{(4.4.6) est \unsure{prouvée}.}
\deleted{à g. de $TC$ (4.4.4) et du fait que $\unsure{C_0}\subset\mathrm{Cof}_W$}
L'équnce de a), b), c) n'est autre que
la prop. 4.3, appliquée aux $W$ et aux deux
couples de Quillen $(C,TF),(\widetilde{TC_0},F_0)$ (compte
tenu du th. 1 du \S\,1). \deleted{et} Qu'on a donc
une \unsure{structure|str} de Quillen close a été vu
dans le cor. 3 de prop. 4.3 (p. 35). \deleted{\unsure{Comme}}
Il reste à prouver l'équivalence de
a et de a'. On a $\mathrm a\Rightarrow\mathrm a'$, puisque
$(TC_0)_*=(\widetilde{TC_0})_*=F_0$. \deleted{Mais} Inversement,
a') implique $\widetilde{TC_0}=\widetilde{TC}$, donc le th. de
factorisation pour $(TC,F)$ en vertu du
th. de factorisation cor. 1 (p. 11) appliqué à
\deleted{couple} $TC_0\subset TC\subset\widetilde{TC_0}$ et à $F_0=F=(TC_0)_*$.

% archive sheet number: 167
